% pro tvar "isosceles triangle" a anchor .apex + výpočty souřadnic
\usetikzlibrary{shapes.geometric,calc}

\begin{tikzpicture}[
    scale=1.35,
    every node/.style={
        circle,
        draw,
        fill=white,
        inner sep=0pt,
        minimum size=7pt
    },
    line width=0.7pt
]

% ============================================================
% GLOBÁLNÍ PARAMETRY
% ============================================================

% Vzdálenost sourozeneckých vrcholů
\def\SiblingDist{0.65}

% Vertikální vzdálenost mezi úrovněmi
\def\LevelDist{0.60}

\def\SubtreeLabelOffset{0.28} % svislá mezera mezi trojúhelníkem a popiskem


% ============================================================
% A_1
% ============================================================

\node (A1) at (0,0) {};

\node[draw=none, fill=none, circle] at (0,0.45)
    {$A_1=1$};


% ============================================================
% A_2
% ============================================================

\node (A2root) at (1.7,0) {};

\node (A2child) at
    ({1.7-\SiblingDist/2},-\LevelDist) {};

\draw (A2root) -- (A2child);

\node[draw=none, fill=none, circle] at (1.7,0.45)
    {$A_2=2$};


% ============================================================
% A_3
% ============================================================

\node (A3root) at (3.4,0) {};

\node (A3left) at
    ({3.4-\SiblingDist/2},-\LevelDist) {};

\node (A3right) at
    ({3.4+\SiblingDist/2},-\LevelDist) {};

\node (A3leftchild) at
    ({3.4-\SiblingDist},-{2*\LevelDist}) {};

\draw (A3root) -- (A3left);
\draw (A3root) -- (A3right);
\draw (A3left) -- (A3leftchild);

\node[draw=none, fill=none, circle] at (3.4,0.45)
    {$A_3=4$};


% ============================================================
% A_4
% ============================================================

\node (A4root) at (5.35,0) {};

% druhá úroveň
\node (A4left) at
    ({5.35-\SiblingDist/2},-\LevelDist) {};

\node (A4right) at
    ({5.35+\SiblingDist/2},-\LevelDist) {};

% třetí úroveň
\node (A4leftleft) at
    ({5.35-\SiblingDist},-{2*\LevelDist}) {};

\node (A4leftright) at
    (5.35,-{2*\LevelDist}) {};

\node (A4rightright) at
    ({5.35+\SiblingDist},-{2*\LevelDist}) {};

% čtvrtá úroveň
\node (A4bottom) at
    ({5.35-1.5*\SiblingDist},-{3*\LevelDist}) {};

% hrany
\draw (A4root) -- (A4left);
\draw (A4root) -- (A4right);

\draw (A4left) -- (A4leftleft);
\draw (A4left) -- (A4leftright);

\draw (A4right) -- (A4rightright);

\draw (A4leftleft) -- (A4bottom);

\node[draw=none, fill=none, circle] at (5.35,0.45)
    {$A_4=7$};


% ============================================================
% Rekurence A_h
% ============================================================

\node[
    draw=none,
    fill=none,
    circle,
    anchor=west
] at (6.55,0.45)
    {$A_h=A_{h-1}+A_{h-2}+1$};


% ============================================================
% kořen
% ============================================================

\node (Ahroot) at (7.9,0) {};


% ============================================================
% levý podstrom h-1
% ============================================================

\node[
    isosceles triangle,
    shape border rotate=90, % špička nahoru (místo doprava)
    anchor=apex,            % zarovnávej podle horní špičky
    draw,
    fill=white,
    minimum height=44pt,
    minimum width=26pt,
    inner sep=0pt
] (Ahleft) at
    ({7.9-\SiblingDist},-{1.25*\LevelDist})
    {};

% hrana do špičky trojúhelníku
\draw (Ahroot) -- (Ahleft.apex);


% ============================================================
% pravý podstrom h-2
% ============================================================

\node[
    isosceles triangle,
    shape border rotate=90, % špička nahoru (místo doprava)
    anchor=apex,            % zarovnávej podle horní špičky
    draw,
    fill=white,
    minimum height=26pt,
    minimum width=16pt,
    inner sep=0pt
] (Ahright) at
    ({7.9+\SiblingDist},-{1.25*\LevelDist})
    {};

% hrana do špičky trojúhelníku
\draw (Ahroot) -- (Ahright.apex);


% ============================================================
% popisky podstromů
% ============================================================

% umísti popisky vždy pod trojúhelník (nezávisle na jeho velikosti)
\node[draw=none, fill=none] at ($(Ahleft.south)+(0,-\SubtreeLabelOffset)$) {$h-1$};
\node[draw=none, fill=none] at ($(Ahright.south)+(0,-\SubtreeLabelOffset)$) {$h-2$};
    
\end{tikzpicture}